% Diagram datových struktur překladače (dle src/*.h na větvi main).
% Každý typ je v diagramu právě jednou; rekurzivní struktury (strom scopů,
% AST, uzly BST) jsou znázorněny smyčkou na sebe sama.
% Soubor je samostatný: styly jsou definované přímo v tikzpicture, vyžaduje
% pouze balík tikz a knihovnu arrows.meta (obojí je v preambuli main.tex).
% čárkovaná šipka = vypůjčený ukazatel (bez vlastnictví).

\begin{figure}[H]
  \centering
  \begin{tikzpicture}[
      ds/.style={draw, rounded corners=1mm, fill=black!4, align=left,
                 inner sep=4pt, font=\footnotesize},
      own/.style={-{Stealth[length=2mm]}, semithick},
      ref/.style={-{Stealth[length=2mm]}, semithick, dashed, black!60},
      lab/.style={font=\scriptsize\ttfamily, fill=white, inner sep=1pt},
    ]
    % ---------------------------------------------------------------- uzly
    \node[ds] (prog) at (0,0)
      {\textbf{Program}\\
       \texttt{functions, initialized}};
    \node[ds] (ft) at (0,-1.8)
      {\textbf{FunctionTable} {\scriptsize(dynamické pole)}\\
       \texttt{functions[], function\_count}};
    \node[ds] (fs) at (0,-4.0)
      {\textbf{FunctionSymbol} {\scriptsize(jedna funkce)}\\
       \texttt{name, return\_type, argument\_count}\\
       \texttt{parameters[]}$^{*}$\texttt{, local\_scope, ast\_root}\\
       \texttt{callees[], is\_main, is\_builtin, \dots}};

    \node[ds, anchor=north] (st) at (-3.0,-6.0)
      {\textbf{SymbolTable} {\scriptsize(jeden scope)}\\
       \texttt{parent, children[], bindings}};
    \node[ds] (bst) at (-3.0,-8.6)
      {\textbf{BST} {\scriptsize(výškově vyvážený)}\\
       \texttt{root, compare\_data, \dots}};
    \node[ds] (bnode) at (-3.0,-10.8)
      {\textbf{Node} {\scriptsize(uzel stromu)}\\
       \texttt{value, height, left, right}};
    \node[ds] (sym) at (0.6,-12.9)
      {\textbf{Symbol} {\scriptsize(proměnná / parametr)}\\
       \texttt{name, type, initialized,}\\
       \texttt{constant, owner, assembler\_suffix}};

    \node[ds, anchor=north] (ast) at (3.3,-6.0)
      {\textbf{ASTNode} {\scriptsize(příkaz / výraz)}\\
       \texttt{kind, inferred\_type,}\\
       \texttt{line, column,}\\
       \texttt{data} -- union podle \texttt{kind}:\\
       \texttt{block, if\_statement,}\\
       \texttt{while\_loop, for\_loop,}\\
       \texttt{assignment, call,}\\
       \texttt{binary\_op, unary\_op, \dots}};

    % ---------------------------------------------------------- vlastnictví
    \draw[own] (prog) -- node[lab, right] {functions} (ft);
    \draw[own] (ft) -- node[lab, right] {functions[i] \normalfont\scriptsize(1\,:\,n)} (fs);
    \draw[own] ([xshift=0.9cm]fs.south west) -- node[lab, above left]
      {local\_scope \normalfont\scriptsize(kořenový scope)} (st.north);
    \draw[own] ([xshift=1.7cm]fs.south) -- node[lab, above right]
      {ast\_root \normalfont\scriptsize(kořen: \texttt{AST\_BLOCK})} (ast.north);

    % strom scopů: vnořené scopy jsou tentýž typ SymbolTable
    \draw[own] ([yshift=0.2cm]st.west) .. controls +(150:1.2) and +(210:1.2) ..
      node[lab, left, align=right] {children[i]\\\normalfont\scriptsize(0\,:\,n)}
      ([yshift=0.15cm]st.south west);
    \draw[own] (st) -- node[lab, left] {bindings} (bst);
    \draw[own] (bst) -- node[lab, left] {root} (bnode);
    \draw[own] ([yshift=0.2cm]bnode.west) .. controls +(150:1.2) and +(210:1.2) ..
      node[lab, left] {left, right} ([yshift=0.15cm]bnode.south west);
    \draw[own] (bnode.south) to[out=-90, in=180]
      node[lab, pos=0.4, below left] {value} (sym.west);

    % AST: potomci uzlu jsou opět uzly ASTNode
    \draw[own] ([yshift=0.6cm]ast.east) .. controls +(20:1.6) and +(340:1.6) ..
      node[lab, right, align=left] {\normalfont\scriptsize potomci\\(0\,:\,n)}
      ([yshift=0.6cm]ast.south east);

    % ------------------------------------------------- vypůjčené ukazatele
    \draw[ref] ([yshift=0.15cm]st.south east) .. controls +(330:0.9) and +(30:0.9) ..
      node[lab, right] {parent} ([yshift=0.2cm]st.east);
    \draw[ref] ([xshift=0.9cm]ast.south west) to[out=-90, in=0]
      node[lab, pos=0.45, right] {variable.symbol, target} (sym.east);
    \draw[ref] ([xshift=1.4cm]ast.north) to[out=90, in=0]
      node[lab, pos=0.6, right] {call.function} (fs.east);
  \end{tikzpicture}
  \caption{Diagram datových struktur překladače. Každý typ je uveden jednou,
    smyčka značí rekurzivní strukturu: tabulky symbolů tvoří strom
    zanořených scopů (funkce odkazuje pouze na kořenový scope, vnořené
    scopy jsou jeho potomci), uzly AST mají za potomky opět uzly AST (např.\ příkazy bloku, operandy
    výrazu, větve podmínky nebo tělo cyklu).
    Plná šipka značí vlastnictví (podstruktura je uvolněna spolu
    s~rodičem), čárkovaná šipka vypůjčený ukazatel bez vlastnictví.
    $^{*}$Parametry jsou vypůjčené symboly z~kořenového scope funkce.}
  \label{fig:datove-struktury}
\end{figure}
